\documentclass[main.tex]{subfiles}
\begin{document}
\section{Discussion and Future Work}
\label{sec:discussion}

We initially began this project by attempting to develop the optimizer directly,
  but discovered that it was far too difficult to write correct rewrite rules:
  as we have discussed in the paper, most require various side conditions,
  and developing the side conditions requires careful reasoning
  about both geometry and the semantics of Skia,
  which at the time we didn’t have a good grasp of.
An example of this is the complexity involved in dealing with \Clip{s}.
The \Clip command,
  together with the \Transform command described in \Cref{subsec:lean},
  are imperative in nature:
  they modify, instead of replacing, the clip (and transform) state.
So, we were ``forced'' to make the \Layer language to make this clip (and transform) state explicit and functional.
Thus we changed our focus to develop the formalization,
  which then enabled us to rapidly develop the rewrite rules (described in \Cref{sec:rewrites}) and
  the optimizer (described in \Cref{sec:compiler}).

Our experience suggests several directions for future work on improving
  the formalization and optimization of Skia.

Skia supports many representations of shapes,
  including rectangles, pixel-boundary-aligned rectangles,
  rounded rectangles, polygons, paths, and text.
Although paths provide a general representation, Skia often handles more specific shape types through specialized code paths.
Transformations between these representations can also be confusing;
for example, the intersection of two rectangles is a rectangle,
  but the intersection of two rounded rectangles is a path.
Currently, our abstract model coalesces all these shapes into various constructors of one type, \Shape.
Future work could track the different types of shapes and the allowable transformations between them,
  which could enable more granular rewrite rules.

Mask filters, which are commonly used to apply Gaussian blurs,
  interact with shape boundaries in ways that are not yet fully captured by our model.
Future work could formalize these interactions.
This direction may be particularly relevant as
  the Skia team has acknowledged the complexity of mask filters
  and are trying to replace them with image filters to simplify the API.

In the course of exploring Skia we noticed that Skia and other 2D graphics libraries share a lot of features,
  and often have quite similar APIs too.
Hence, one direction we plan to explore in the future
  is expanding our formalization to other graphics APIs so that,
  for example, we could verify a compiler from, say, SVG or HTML <canvas> to Skia or another back-end library.
We also hope to explore compositing,
 a type of caching used in animations,
 and how it impacts graphics library correctness and performance.

\end{document}
